% Page 020

\seite{20}

\jahresueberschrift{Das Jahr 1892}
\pstart
Die Schülerzahl beträgt in diesem Jahre 50.\ob
Im Winter dieses Jahres geriet der Gemeinderat mit dem\ob
größten Teile der Bürger in Konflikt, weil derselbe den\ob
Entschluß faßte, ein Drittel des Waldbesitz den Berechtigten\ob
zu entziehen. Von den 46 Berechtigten des Vorrats berg\ob
jeder grundsätzlich schätze 6 Raumeter Buchenholz\ob
\margmark{über Waldmeiste\\Klara Ortschaft}%
nicht besser und zwei Raumeter Tannenholz. 32\ob
Bürger legten Berufung ein bei dem zuständigen\ob
Kreisausschusse, jedoch ohne jeglichen Erfolg. Das Geld dazu zur\ob
Aufsteigerung.

Die Gemeinde Sefferweich hat ungefähr 300 Morgen Wald,\ob
davon sind 300 Morgen Kiefer und 500 M.\ Gesamtwert.\ob
Sefferweich liegt auf dem einzigen nächsten\ob
Siebengemeinden Wald. In früherer Zeit konnte die\ob
Gemeinde sämtliche Offerten im Offizielt Klaftbrunnen\ob
erhalten, 1500 Morgen, wovon sie auf ihren Antheil\ob
an Siebengemeinden Wald verzichtete. Die Gemeinde\ob
\margmark{Die früheren Bür-\\gergüter auf dem\\Klaftbrunnen.}%
ging nicht darauf ein, weil es nach dem Urteil\ob
von Sachverständigen ein nachteilhaftes Tauschgeschäft ange\ohb
nommen.

Auf dem Klaftbrunnen befindet sich ferner ein kleines\ob
Kapelle, die jetzt in einen Verfall zugewandelt ist. Über\ob
die Kapelle war die Abstimmung des Herrn Pfarrer Dunckel, die\ob
für alljährlich bei St.\ Markus Lud. Die Bausubstanz der Kapelle ist\ob
Seite noch aufrecht. (Derzeitiger Vorsteher: Pitts Johann)
\pend
